%! AP Statistics — Unit 3, Lesson 3.7 — Warm-Up
\documentclass[10pt]{article}
\usepackage{apstats-article}
\usepackage{apstats-boxes}

\begin{document}

\pageheader{Unit 3, Lesson 3.7}{Warm-Up}
\namedateperiod

\vspace{0.3in}

% ── SECTION 1: Spiral Review ──────────────────────────────────────────────────
{\large\bfseries\color{navy} Part 1 \quad Spiral Review}

\vspace{0.12in}

\begin{enumerate}

    \item From Lessons 3.5--3.6 --- for each experiment, recommend a design and
    justify your choice.

    \vspace{8pt}
    \renewcommand{\arraystretch}{1.8}
    \begin{tabularx}{\linewidth}{@{} X p{2.8cm} p{3.4cm} @{}}
    \textbf{Scenario} & \textbf{Design} & \textbf{Brief justification} \\
    \hline
    Test a new fertilizer on 30 identical corn plants. &
    \blank{2.5cm} & \blank{3.2cm} \\
    Test two teaching methods on 40 students who vary widely in reading ability. &
    \blank{2.5cm} & \blank{3.2cm} \\
    \end{tabularx}

    \vspace{0.18in}

    \item Two studies find the same result --- patients taking the new drug recover
    faster. In Study A, 50 volunteers chose which treatment to take. In Study B,
    50 randomly selected patients were randomly assigned to drug or placebo.
    \begin{enumerate}[label=(\alph*), itemsep=4pt]
        \item Which study can conclude the drug \emph{caused} faster recovery?
        \blank{3cm}
        \item Which study's results can be generalized to the broader patient
        population? \blank{3cm}
        \item What is different about the two studies that leads to these different
        conclusions?\\[4pt]
        \writeline\writeline
    \end{enumerate}

\end{enumerate}

\vspace{0.2in}

% ── SECTION 2: Looking Ahead ──────────────────────────────────────────────────
{\large\bfseries\color{navy} Part 2 \quad Looking Ahead}

\vspace{0.12in}

\noindent Today we build the \textbf{scope of inference} --- the complete framework for
what we can conclude and about whom.

\vspace{0.14in}

\begin{tcolorbox}[colback=greenbg, colframe=greenacc, boxrule=0.8pt, arc=2mm,
    left=4mm, right=4mm, top=3mm, bottom=3mm]
    \small
    \begin{itemize}[leftmargin=1.3em, itemsep=8pt]
  \item Random \textbf{assignment} $\Rightarrow$ can conclude \emph{causation}.
  \item Random \textbf{selection} $\Rightarrow$ can \emph{generalize} to the population.
  \item Both $\Rightarrow$ can conclude causation and generalize.
  \item Neither $\Rightarrow$ can only describe the sample.
    \end{itemize}
\end{tcolorbox}

\vspace{0.2in}

\noindent\textcolor{slate}{\small In your own words: why does random assignment allow
causal conclusions but random selection does not?}\\[8pt]
\writeline\\[2pt]\writeline

\end{document}
